\paragraph{Data and forecasts.}
We use daily closing prices of seven equity indices (S\&P 500, FTSE 100,
DAX, Nikkei 225, Euro Stoxx 50, Hang Seng and Bovespa) from January 2000 to
September 2026, and percentage log returns. Six standard models produce
one-day-ahead (VaR, ES) forecasts at $\tau=0.025$ and $\tau=0.01$, each
using returns up to the previous day: historical simulation on the last 250
and 500 returns (HS-250, HS-500); the exponentially weighted variance with
decay factor 0.94 and normal quantiles (EWMA); GARCH(1,1) with normal (GARCH-N)
and standardized Student-$t$ innovations (GARCH-t); and filtered historical
simulation (FHS), which multiplies the GARCH-N volatility by the empirical
quantile and tail mean of the standardized residuals. The GARCH models have
zero mean and are re-estimated every 20 days on the last 1000 returns. After
an estimation window of 1000 days and a further 250 days that fix the
scale $\lambda_0$ of the untuned confidence sequence, the evaluation
periods have between 3{,}638 and 5{,}541 days.

\paragraph{Procedures.}
For each index and level we compare the 15 pairs of models by the FZ0 loss.
The Diebold--Mariano test is two-sided at 5\% with a Newey--West variance
and lag $\lfloor4(T/100)^{2/9}\rfloor$. We apply it to the full evaluation
period and to each non-overlapping window of 250 days, the length of the
regulatory backtesting sample. In each window we record whether a
significant result favours the model whose ES forecast was less
conservative on average over the window, and whether its sign agrees with
that of the mean loss difference over all other days of the evaluation
period. The certification of Proposition~\ref{prop:es-one-sided} is applied
on the days on which each model is the less conservative, with the outcome
floored at $\underline y_t=\kappa\min\{e_{A,t},e_{B,t}\}$ and $\kappa=2$,
that is, at twice the more conservative ES forecast; $\delta=0.05$, $K=20$,
and on the full period the sequence is tuned on the first 10\% of the
selected days. A tuning segment of 10\% of a 250-day window would contain
less than one expected exceedance, so within windows we use the untuned
sequence, with $\lambda_0$ from the 250 days that precede the evaluation
period. The choice of $\kappa$ affects power but not validity; we also
report $\kappa=1.5$ and $\kappa=3$.

\paragraph{Results.}
Table~\ref{tab:es} summarizes the results. A 250-day window contains a
median of nine days on which the return falls below at least one of the two
VaR forecasts at $\tau=0.025$, and four at $\tau=0.01$, against about six
and 2.5 for one calibrated forecast, because two forecasts are involved and
the models' exceedance rates, pooled over indices, are 1.0--1.4 times the
nominal level at $\tau=0.025$ and 1.1--2.1 times at $\tau=0.01$. The Diebold--Mariano test
is significant in 18\% of the window--pair tests at both levels, and 73\%
(at $\tau=0.025$) and 80\% (at $\tau=0.01$) of the significant results
favour the less conservative model. Over the full period, with a median of
203 and 104 exceedances per pair, the corresponding shares are 67\% and
51\%, counting a model as less conservative when its ES was the less
conservative on most days. The tilt towards the less conservative model is
thus modest at $\tau=0.025$, 73\% in windows against 67\% over the full
period, and large at $\tau=0.01$, 80\% against 51\%, where windows hold
the fewest exceedances. The significant window results that favour the less
conservative model contradict the ranking on the other days in 18.4\% and
21.1\% of cases; those that favour the more conservative model do so in
6.9\% and 10.5\%.

The window tests are dependent: the 15 pairs of an index share the returns,
and windows of different indices overlap in time. To assess the sampling
uncertainty of these shares, we assign each window to the calendar year of
its midpoint and resample years with replacement, keeping every window of
the drawn years (4000 resamples). The 95\% percentile interval for the
difference between the two wrong-sign rates is $[1.7,19.7]$ percentage
points at $\tau=0.025$ and $[0.0,18.9]$ at $\tau=0.01$; for the share of
significant results that favour the less conservative model it is
$[54,89]\%$ and $[67,90]\%$. The asymmetry is thus clear at
$\tau=0.025$ and borderline at $\tau=0.01$, where significant results of
either kind are fewer.

Replacing the normal critical values of the Diebold--Mariano test by those
of a circular block bootstrap-$t$ test (blocks of 10 days, 999 resamples,
with the statistic and its replicates studentized by the variance of the
block means) removes most of the tilt. The share of significant window
tests is similar, 18\% and 20\%, but only 51\% and 60\% of them favour the
less conservative model. Their wrong-sign rates are 12.9\% and 14.2\%,
against 9.5\% and 6.3\% for those that favour the more conservative
model, so at $\tau=0.01$ results
in favour of the less conservative model remain the less reliable.

The certification never succeeds within a window, as the lower bound of
Section~\ref{sec:es} leads one to expect on four to nine exceedances. On
the full period it certifies that the less conservative model was worse
on its own days for 48\% of pairs at $\tau=0.025$ and 44\% at
$\tau=0.01$. Across $\kappa\in\{1.5,2,3\}$ these shares range over 48--50\%
and 42--44\%. The untuned sequence certifies 43\% and 24\%. When it certifies, the
full-period
Diebold--Mariano test is significant in the same direction in 48 of 50 and
36 of 46 pairs, and never in the opposite direction.

\begin{table}[t]
\centering\small
\caption{Comparisons of six (VaR, ES) models on seven equity indices,
15 pairs each, by the FZ0 loss. Windows: non-overlapping 250-day windows.
Exceedances: days on which the return falls below the larger VaR forecast.
DM: two-sided 5\% Diebold--Mariano test with normal critical values. Wrong
sign: share of the significant window results, of each kind, whose sign is
opposite to that of the mean over all other days. EB certifies: share for
which the one-sided 95\% bound of Proposition~\ref{prop:es-one-sided}
certifies that the less conservative model was worse on its own days, with
the outcome floored at twice the more conservative ES; untuned sequence in
windows, tuned sequence over the full period.}
\label{tab:es}
\begin{tabular}{lcc}
\toprule
 & $\tau=0.025$ & $\tau=0.01$\\
\midrule
\multicolumn{3}{l}{\emph{250-day windows}}\\
window--pair tests & 2130 & 2130\\
median exceedances & 9 & 4\\
DM significant (\%) & 17.8 & 18.0\\
\quad favouring the less conservative (\%) & 73.1 & 80.2\\
\quad wrong sign, favouring less conservative (\%) & 18.4 & 21.1\\
\quad wrong sign, favouring more conservative (\%) & 6.9 & 10.5\\
EB certifies (\%) & 0.0 & 0.0\\
\midrule
\multicolumn{3}{l}{\emph{full period}}\\
pairs & 105 & 105\\
median exceedances & 203 & 104\\
DM significant (\%) & 77.1 & 71.4\\
\quad favouring the less conservative (\%) & 66.7 & 50.7\\
EB certifies (\%) & 47.6 & 43.8\\
\bottomrule
\end{tabular}

\end{table}
